\section{Discussion and Limitations}

WhiteTiger provides a standardized and temporally structured data foundation for embodied intelligence research at billion-frame scale. Its principal properties are fully human-teleoperated physical collection, 14-platform embodiment coverage, preservation of platform-specific observation and action schemas, LeRobot v2.1 conversion, and event-level temporal annotations. With \textbf{12877.59 hours}, \textbf{737575 episodes}, and more than \textbf{1.39 billion frames}, the corpus supports foundation-model adaptation regimes that are difficult to study with smaller task-specific datasets.

An offline open-loop zero-shot evaluation is planned to assess action-prediction transfer from WhiteTiger fine-tuning to external robot demonstration datasets excluded from WhiteTiger training. The planned protocol will compare a pretrained baseline with a WhiteTiger-fine-tuned model on identical held-out trajectories without adaptation to the evaluation datasets. Because the experiment has not yet been completed, the current version does not report quantitative transfer results or claim that WhiteTiger fine-tuning improves offline prediction performance.

WhiteTiger event annotations provide temporal structure beyond dense action labels. They expose task progress, subgoal completion, contact transitions, alignment states, boundary crossings, deformation states, and object-state changes that are often implicit in raw trajectories. This structure supports keyframe-aware learning paradigms, including waypoint-conditioned imitation learning, goal-conditioned policy learning, hierarchical skill decomposition, reward shaping, value-guided policy improvement, subgoal-image reasoning, and phase-aware evaluation. The annotations should be interpreted as a general temporal supervision resource rather than labels optimized for a single algorithm.

Offline open-loop action prediction will not provide a complete measure of real-world robot performance. Normalized prediction error under recorded observations can measure agreement with reference teleoperation actions, but it cannot evaluate closed-loop recovery, contact-rich interaction stability, or task success under policy-induced state distributions. Sequential decision making can compound prediction errors as future observations shift away from demonstration trajectories~\cite{ross2011dagger,zhao2023aloha}. The planned benchmark should therefore be interpreted as a transfer-oriented diagnostic rather than evidence of real-world task success or held-out-platform generalization.

The initial planned benchmark uses one model family and one training comparison. Additional baselines such as ACT, Behavior Transformer, Diffusion Policy, and open-source vision-language-action policies would broaden comparative coverage~\cite{zhao2023aloha,shafiullah2022behaviortransformer,chi2023diffusionpolicy,kim2024openvla}. Data-scaling studies, held-out-platform evaluation, closed-loop robot rollouts, trajectory-weighted aggregation, and uncertainty estimates across evaluation trajectories are also required to characterize which dataset factors govern policy scaling and generalization.

Event- and keyframe-level annotations remain a dataset capability and a planned experimental direction in the current version. The real-world evaluation of event-guided value learning has not yet been completed. Claims about improved real-robot execution from event-derived supervision should therefore be deferred until policies trained with and without event-based value signals are compared under identical closed-loop tasks with execution success rate, event completion rate, failure-stage distribution, and time-to-completion metrics.
